\documentclass[12pt]{article}
\usepackage[utf8]{inputenc}
\usepackage[T1]{fontenc}
\usepackage{hyperref}
\title{Wusa Nasu Language Resources for K-12 Curriculum Development}
\author{Compiled by Claude Sonnet 5.5 (claude-sonnet-5-5) for Andr\'as Kornai}
\date{}
\begin{document}
\maketitle

\section{Introduction}
Wusa Nasu (yig) is an Eastern Yi (Tibeto-Burman) language of northwest Guizhou and northeast Yunnan. It is strongest in Weining, Shuicheng, Hezhang and Nayong counties and in Xuanwei. Population estimates vary widely, from about 250,000 (a mission profile) to 800,000 (GRN). One source reports 85\% of Yi speakers in Weining using their native tongue daily. A distinct form of the Yi script was traditionally used but few can still read it. Search found no textbooks or school programmes.

\section{Primary Educational Resources}
None found.

\section{Linguistic and Reference Materials}
\begin{itemize}
\item Nasu language (Wikipedia): \url{https://en.wikipedia.org/wiki/Nasu_language}
\item Nasu people (Wikipedia): \url{https://en.wikipedia.org/wiki/Nasu_people}
\end{itemize}

\section{Religious and Community Resources}
\begin{itemize}
\item Asia Harvest profile of the Wusa Nasu (a mission profile): \url{https://people-groups.asiaharvest.org/people-group-profiles/guizhou/nasu-wusa.pdf}
\end{itemize}

\section{Digital and Audio Resources}
\begin{itemize}
\item Global Recordings Network page for Nasu, Wusa: \url{https://www.globalrecordings.net/en/language/yig}
\end{itemize}

\section{Curriculum Development Context}
The language is widely spoken, but no native-language educational material was found. A curriculum would have to be developed from scratch. Sources disagree on the varieties covered by the label ``Nasu'' and on the script, so the target variety and orthography would need to be settled first.
\end{document}